\chapter{绪论}
\label{chap:introduction}

\section{研究背景与意义}

\subsection{检索增强生成的发展现状}

大语言模型（Large Language Models, LLMs）的快速发展是近年来人工智能领域最引人注目的技术突破之一。以 GPT-4~\citep{openai2023gpt4}、Gemini~\citep{geminiteam2024gemini15}、DeepSeek-V3~\citep{deepseekai2024deepseekv3}、Qwen2.5~\citep{qwen2025qwen25technicalreport} 为代表的大语言模型，在文本理解、逻辑推理、代码生成等广泛的自然语言处理任务上展现出了令人瞩目的能力，深刻改变了人机交互的范式。然而，尽管大语言模型的规模不断扩大、能力持续增强，其固有的局限性也日益凸显。

首先，知识时效性受限。大语言模型的知识来源于预训练阶段所接触的静态语料库，存在明确的"知识截止日期"，无法获取训练数据截止日期之后的新信息和新事件。其次，事实性幻觉问题严重。研究表明，大语言模型在生成过程中倾向于产生看似流畅合理但实际上与客观事实不符的内容~\citep{xu2025hallucination}，这种现象在医疗、法律、金融等对准确性要求极高的领域构成了严重的安全风险。第三，长尾知识与领域知识匮乏。Kandpal 等人的研究指出，大语言模型在学习长尾知识方面存在显著困难~\citep{kandpal2023largelanguagemodelsstruggle}——对于训练语料中出现频率较低的实体和事实，模型的记忆与生成能力急剧下降。此外，对于特定垂直领域或企业私有知识库中的专业知识，通用大语言模型往往缺乏足够的覆盖。

为了克服上述局限性，检索增强生成（Retrieval-Augmented Generation, RAG）技术应运而生，并迅速成为当前大语言模型应用中最核心的技术范式之一~\citep{gao2024retrievalaugmented}。RAG 的核心思想是在模型生成回答之前，首先从外部知识库中检索出与用户查询相关的文档或段落，然后将这些检索到的证据作为上下文提供给语言模型，从而引导其生成更加准确、有据可查的回复。这种将参数化记忆（模型内部权重）与非参数化记忆（外部知识库）相结合的思路，使得模型可以动态地访问最新信息，有效缓解了知识过时和幻觉问题~\citep{shuster2021retrieval}。Lewis 等人~\citep{lewis2020rag} 与 Guu 等人~\citep{guu2020realm} 的里程碑式工作奠定了现代 RAG 研究的基本范式，此后 RAG 已从辅助手段成长为连接大模型能力与行业知识的核心基础设施。

从技术架构来看，现代 RAG 系统已经从最初的"检索-拼接-生成"朴素流水线，演化为包含查询理解、多路检索、重排序、上下文压缩、自适应检索等多个精细化模块的复杂系统。其中，\textbf{检索模块的召回能力决定了整个系统的上限}——即使生成端拥有再强的语言模型，若检索到的证据不相关或不完整，后续生成无法无中生有。以 DPR~\citep{karpukhin-etal-2020-dense}、Contriever~\citep{izacard2022contriever}、BGE~\citep{xiao2023cpack} 为代表的密集向量检索模型取代了传统的稀疏检索（如 BM25），在单轮场景下已具备高度成熟的语义匹配能力。

\subsection{多轮对话检索与查询重写的核心挑战}

随着 ChatGPT 等对话式产品的普及，用户与 RAG 系统的交互方式已从单轮独立查询向\textbf{多轮对话}形态快速迁移。在真实场景下，用户往往不会一次性把所有信息表达清楚，而是通过连续多轮追问、澄清、切换话题来逐步完成信息获取。这种\textit{多轮对话检索}（Multi-turn Conversational Retrieval）场景给系统带来了远比单轮更为复杂的挑战。

考察以下一个真实多轮对话片段：

\begin{quote}
\textbf{U1}：What is the capital of France? \\
\textbf{A1}：Paris. \\
\textbf{U2}：Who founded it? \\
\textbf{A2}：Paris was founded by the Parisii tribe in the 3rd century BC. \\
\textbf{U3}：And its population now?
\end{quote}

当检索系统需要为 U3 检索相关证据时，直接把 "And its population now?" 输入检索器将无法得到任何有意义的结果——"its" 指代不明，句子的主语被省略，信息完全依赖对话历史。只有把查询补全为"What is the current population of Paris?" 这样的自包含查询，现代稠密检索器才能返回正确的证据文档。这种\textbf{把上下文依赖的用户问句改写为语义自包含的独立查询}的技术，即为\textbf{查询重写}（Query Rewriting, QR），是多轮对话 RAG 系统中承上启下的关键环节~\citep{mo2024convsar}。

真实多轮对话中带来查询语义不完整的语言现象主要包括：
\begin{itemize}
\item \textbf{指代（Coreference）}：使用"他/她/它/前者/这个"等代词替代先前提到的实体~\citep{lee2017coref}；
\item \textbf{省略（Ellipsis）}：省略主语、谓语、宾语或修饰成分，依赖上下文补全~\citep{halliday1976cohesion,quan2019gecor}；
\item \textbf{话题漂移（Topic Drift）}：后续轮次可能与早期轮次话题不再直接相关，甚至发生硬切换~\citep{adlakha-etal-2022-topiocqa}；
\item \textbf{隐式澄清}：用户基于系统前一轮回答进行追问，查询的含义依赖于答案的内容。
\end{itemize}

针对上述挑战，业界主要存在两条技术路线。一条是\textit{端到端对话检索器}（如 Dragon-ChatQA~\citep{Liu2024ChatQASG}），直接把对话历史与当前轮次拼接作为检索输入，让模型自行学会消解上下文依赖；另一条则是\textit{查询重写}~\citep{elgohary2019canyouunpack,vakulenko2021question,mao2023llm4cs}，在检索之前显式地把当前轮次改写为独立查询，再交给任意现成的单轮检索器。本文聚焦于后者，原因在于：（1）查询重写与下游检索器完全解耦，可即插即用地与任意单轮检索器组合，部署灵活性显著高于专用对话检索器；（2）重写后的查询可直接复用于排序、摘要、用户意图分析等多个下游模块，具有更强的工程复用价值；（3）现有对话检索器的训练依赖昂贵的对话标注数据，而查询重写方法可以充分利用日益丰富的大语言模型上下文理解能力，研究空间更大。

\subsection{现有查询重写方法的瓶颈}

尽管围绕查询重写已有一批代表性工作~\citep{voskarides2020quretec,yu2020fewshotgenerative,wu2022conqrr,mao2023llm4cs,mo2023convgqr,jang2024itercqr}，现有方法仍远未解决多轮对话检索的核心难点。本文将其归纳为三大瓶颈：

\textbf{瓶颈一：历史噪声建模不足。}绝大多数已有方法在重写阶段将对话历史 $H_i=\{(q_1,a_1),\dots,(q_{i-1},a_{i-1})\}$ 作为完整上下文输入重写模型。但在长对话、跨话题场景下，远早于当前轮次的历史往往已与当前查询无关，甚至存在硬话题切换——此时继续引入无关历史不仅增大输入长度、推高计算成本，更会系统性误导重写器引入无关实体或属性，造成"重写越长、检索越差"的反直觉现象。少数工作尝试用手工规则或启发式截断，但无法区分"同一话题内的弱相关轮次"与"已切换到新话题的噪声轮次"，缺乏对话题漂移的显式建模。

\textbf{瓶颈二：异质现象耦合训练。}现有方法普遍把指代消解与省略补全混合为一个端到端序列生成任务，让模型在单一损失下同时学习两种异质的语言现象。然而，这两种现象在语言学性质、对检索的影响、错误传播模式上差异显著——指代错误通常表现为错词（wrong entity），省略错误通常表现为缺词（missing slot）。端到端训练难以让模型明确地学习两种结构不同的修复操作，导致训练目标不明确、错误难以定位，并且生成的重写在简单样本上"过度改写"、在困难样本上"改写不足"。

\textbf{瓶颈三：重写目标与检索器偏好失配。}现有的有监督重写方法通常将"人工重写的参考答案"作为目标，优化 BLEU/ROUGE 等文本相似度指标~\citep{papineni2002bleu,lin2004rouge}。然而，\textbf{像人}的重写并不等于\textbf{好被检索}的重写——人类标注者通常遵循"尽量用原文词汇且尽量简洁"的直觉，而现代稠密检索器更偏好具备明确实体、明确限定词、句法完整的长查询。这种目标失配导致即使重写器在 BLEU 上取得了高分，下游检索器召回率的提升却十分有限，甚至出现下降。少数工作尝试用检索反馈微调重写器（如 CONQRR~\citep{wu2022conqrr}、IterCQR~\citep{jang2024itercqr}），但普遍采用策略梯度或迭代伪标签，训练不稳定、算力门槛高，难以在中小规模实验室落地。

综上，构建一套\textit{低噪声、结构化、与下游检索器偏好对齐}的查询重写框架，对提升工业级多轮对话 RAG 系统的整体质量具有重要研究价值与实际意义。这正是本文研究工作所致力于解决的核心问题。

\section{国内外研究现状及存在的问题}

围绕多轮对话检索的查询重写任务，国内外已有近十年的研究积累。本节从"早期的监督式重写"、"基于大模型的零样本重写"以及"基于检索反馈的对齐重写"三条技术主线出发，系统梳理现状并指出共性问题。

\subsection{早期的监督式重写方法}

查询重写任务最早可追溯到 Elgohary 等人~\citep{elgohary2019canyouunpack} 提出的 CANARD 数据集，该工作首次以人工众包方式构建了"原始对话问题—去上下文化问题"配对，并训练了基于 Transformer 的 Seq2Seq 重写模型。QuReTeC~\citep{voskarides2020quretec} 将重写任务建模为对历史词项的二元选择——判断每个历史词是否应附加到当前问题后形成扩展查询，从而回避了自由文本生成的训练难度；该方法在 TREC CAsT 上显著提升了 BM25 检索性能，但受限于扩展式重写无法修改句法结构，面对复杂省略与指代时力不从心。Lin 等人~\citep{lin2020conversational}、Vakulenko 等人~\citep{vakulenko2021question} 进一步基于 T5~\citep{raffel2020t5} 和 GPT-2 等预训练语言模型进行端到端微调，得到流畅的重写输出，成为此后一段时间内该方向的主流做法。Yu 等人~\citep{yu2020fewshotgenerative} 则探索了小样本场景下的重写，通过引入弱监督提升低资源性能。

上述监督式方法的共性局限在于：（1）严重依赖稀缺且高成本的人工重写标注；（2）训练目标以词级相似度为主，与下游检索器偏好存在根本失配；（3）未显式建模历史中的噪声轮次，在长对话、跨话题场景下表现不佳。

\subsection{基于大语言模型的零样本重写}

随着 GPT-3.5、GPT-4 等大语言模型的涌现，利用大模型的零样本或少样本能力直接重写对话查询成为可能。LLM4CS~\citep{mao2023llm4cs} 系统性地研究了大模型在对话搜索中的重写与查询扩展能力，提出多路径思维链生成加聚合的重写框架，在 CAsT 系列基准上显著超越此前的监督式方法。Ye 等人~\citep{ye2023enhancing} 则探索大模型辅助生成信息更丰富的重写查询，强调"信息量"而非"简洁性"作为新的重写目标。Mao 等人~\citep{mao2023search} 提出面向搜索的会话查询编辑框架，通过结构化提示词设计让大模型进行有方向的编辑而非无约束改写。

零样本大模型重写极大降低了标注成本，但代价同样明显：（1）高性能通常需要 70B 以上参数规模的闭源模型支撑，推理成本与延迟对工业部署不友好；（2）大模型 prompt 敏感性高，不同提示下输出差异显著，质量不稳定；（3）大模型仍然对完整对话历史"照单全收"，在长对话场景下易被噪声误导；（4）重写目标仍是"像人写"，未考虑下游检索器的真实偏好。

\subsection{基于检索反馈的对齐重写}

为缓解"重写目标与检索偏好失配"的问题，近年来出现了一批基于检索反馈优化重写器的工作。ConvGQR~\citep{mo2023convgqr} 提出同时生成重写查询和潜在答案两条分支，通过两种分支的互补提升检索召回。CONQRR~\citep{wu2022conqrr} 首次将检索器召回质量作为奖励信号，采用 REINFORCE 类策略梯度对 T5 重写器进行强化学习微调，取得了可观但不稳定的性能增益。IterCQR~\citep{jang2024itercqr} 进一步把检索反馈引入迭代自训练循环，通过多轮伪标签更新提升重写器质量。Kim 等人~\citep{kim2022saving} 则从另一角度分析了稠密检索器在对话检索中存在的"捷径依赖"问题，指出检索器会过度依赖当前轮次词汇而忽视历史信息。

这些工作初步验证了"检索反馈驱动重写"的价值，但仍有明显不足：（1）策略梯度方法（REINFORCE、PPO~\citep{schulman2017ppo}）训练不稳定、对奖励设计敏感；（2）迭代伪标签范式存在误差累积风险；（3）现有方法多基于 T5-base/large 规模的重写器，未探索在更强的数十亿参数大模型上的可行性；（4）缺少对"历史剪枝 + 结构化重写 + 检索反馈"三者协同的系统研究。

\subsection{存在的问题小结}

综合以上分析，当前多轮对话检索中查询重写的研究存在三大共性问题：

\begin{enumerate}
    \item \textbf{历史噪声问题}：现有方法普遍将完整对话历史作为重写器输入，缺少对话题漂移的显式建模，在长对话、跨话题场景下性能显著退化；
    \item \textbf{异质现象耦合问题}：指代消解与省略补全被混合为端到端任务，训练信号混杂、错误难以定位；
    \item \textbf{目标失配问题}："像人写"的重写目标与"被检索器召回"的真实目标之间存在系统性偏差，亟需更稳定、更轻量的检索反馈对齐方法。
\end{enumerate}

这三大问题共同构成了本文的研究出发点。

\section{研究目标与主要贡献}

针对上述三大问题，本文系统研究面向多轮对话检索的查询重写方法，提出一套由\textit{历史剪枝—渐进式重写—检索反馈优化}三阶段构成的查询重写框架，并在多个公开基准上开展充分的实证分析。本文的研究目标与主要贡献概括如下：

\textbf{研究目标一：话题漂移感知的对话历史剪枝。}破解"长历史即高噪声"的困境。本文针对性地引入\textbf{轮次级相关性判别器}，基于 ModernBERT-base 构建三分类头（强相关/弱相关/无关），并以 TopiOCQA 话题切换标注与 QReCC 人工重写的词级对比启发式联合构造 94\,422 条弱监督训练样本。该模块独立部署于重写模型之前，将对话历史压缩为高相关子序列后再送入重写器，显著降低输入噪声与推理成本。在 4\,721 条留出验证集上，分类器取得 80.1\% 的准确率与 0.797 的宏 F1（详见 §~\ref{chap:hpprune}）。

\textbf{研究目标二：两阶段渐进式查询重写。}解耦异质语言现象。本文提出\textbf{渐进式重写框架}（Progressive Rewriter, PRW），将重写任务显式分解为两个串行子任务：第一阶段仅完成指代消解，将 "it/they/this" 等代词替换为对应实体而不改变表层句法；第二阶段在第一阶段输出基础上完成省略补全，恢复缺失的主语、谓语与修饰成分，产出自包含查询。两阶段分别以合成监督数据进行有监督微调，并通过阶段一致性损失防止语义漂移。基于 Qwen2.5-3B-Instruct 与 LoRA（$r{=}16$）的两阶段 SFT 实验表明，PRW 在 TopiOCQA、Doc2Dial 两个话题切换密集、上下文依赖强的基准上，将 Recall@5 相对 No-Rewrite 分别提升 39 和 66 个百分点以上（BGE-large 检索器下，详见 §~\ref{chap:prw}）。

\textbf{研究目标三：检索反馈驱动的偏好对齐。}将重写目标真正对齐到下游检索器。本文提出\textbf{检索反馈重写器}（Retrieval-Feedback Rewriter, RFR）方案。具体地，以温度 0.9、top-$p$ 0.95 从 PRW 采样 8 个候选重写，用 BGE-large 计算每个候选与 gold passage 的余弦相似度作为反馈信号，将最优/最差候选构造为偏好对，应用 Direct Preference Optimization~\citep{rafailov2023direct} 对 PRW 的 LoRA 适配器进行轻量化微调。该方案完全无需额外标注、训练稳定、成本低廉。在 4 基准 $\times$ 3 异构检索器共 12 组设定中，RFR 相对 PRW 取得平均 $+1.10$ Recall@5 / $+1.31$ Recall@20 的稳定增益，且该增益在训练从未见过的 GTE-ModernBERT 与 Dragon-Multiturn 检索器上同样成立。

\textbf{综合贡献可归纳为以下四点：}

\begin{enumerate}
    \item 提出了话题漂移感知的对话历史剪枝方法，首次在轮次级引入三分类显式建模，从"噪声过滤"维度提升对话查询重写上限；
    \item 提出了两阶段渐进式重写框架 PRW，将异质语言现象解耦为可单独优化、可单独评估的串行子任务，显著提升重写的可解释性与稳定性；
    \item 提出了基于 DPO 的检索反馈重写方法 RFR，无需额外标注即可将重写目标对齐到下游检索器偏好，在 3B 级别模型上实现了堪比 72B 大模型的重写质量；
    \item 在 QReCC、QuAC、Doc2Dial、TopiOCQA 四个公开基准与一个实验室合作构造的困难评测集上开展了充分的对比实验、消融分析与案例研究\footnote{本文实验中包含的"MTR-Bench"数据集由本实验室合作工作~\citep{mtrsuite2025} 独立构建，本文仅将其作为外部困难评测集使用，相关构建细节不属于本文贡献。}，验证了所提方法在检索召回、排序质量、推理效率与鲁棒性上的综合优势。
\end{enumerate}

\section{论文组织结构}

本文共分为八章，各章节的组织安排如下：

\textbf{第一章\quad 绪论。}介绍本文研究背景、国内外研究现状、研究目标与主要贡献，并概述论文组织结构。

\textbf{第二章\quad 相关工作与技术基础。}系统梳理与本文密切相关的三类工作：（1）检索增强生成与对话检索的技术演进；（2）查询重写方法的发展脉络，包括监督式、零样本大模型、强化微调三条路线；（3）大模型监督微调与直接偏好优化（DPO）等训练方法基础。

\textbf{第三章\quad 问题定义与整体框架。}形式化定义多轮对话查询重写任务，统一符号与评估指标；在此基础上提出本文的"历史剪枝—渐进式重写—检索反馈优化"三阶段总体框架，并给出各模块间的协作流程与训练-推理一致性分析。

\textbf{第四章\quad 话题漂移感知的对话历史剪枝。}详述轮次级相关性分类器的数据构造、模型架构、训练目标与剪枝策略；报告在 TopiOCQA、QReCC 等数据集上的剪枝质量评估与对下游重写效果的消融影响。

\textbf{第五章\quad 两阶段渐进式查询重写。}详述 PRW 框架的任务解耦动机、两阶段训练数据合成、阶段一致性损失设计、以及基于 Qwen2.5-1.5B/3B/7B 的多规模实现；对比端到端基线与单阶段变体，验证解耦设计的必要性。

\textbf{第六章\quad 检索反馈驱动的偏好对齐。}详述 RFR 方案的偏好对构造策略、DPO 训练目标、奖励温度扫描与训练稳定性分析；对比 REINFORCE、PPO 等强化学习替代方案。

\textbf{第七章\quad 实验与结果分析。}在 QReCC、QuAC、Doc2Dial、TopiOCQA 及实验室合作评测集上开展完整实验：主结果对比、逐模块消融、硬话题切换子集分析、推理效率分析、重写长度分布分析、跨检索器泛化测试、以及案例研究。

\textbf{第八章\quad 总结与展望。}总结本文的研究工作与创新点，讨论当前局限性，并对未来可扩展的研究方向进行展望。
